\documentclass[10pt,a4paper]{amsart}
\usepackage[margin=19mm]{geometry}
\usepackage{amsmath,amssymb,mathtools}
\usepackage[scr=rsfs,cal=euler]{mathalfa}
\usepackage{xfrac}
\usepackage[unicode,hidelinks,bookmarks=false]{hyperref}
\newcommand{\ie}{i.e.\ }
\newcommand{\cf}{cf.\ }
\newcommand{\resp}{respectively\ }
\newcommand{\Subsection}{\subsection}
\providecommand{\nobreakdash}{\nobreak\hbox{-}}
\setlength{\emergencystretch}{2em}
\multlinegap0pt
\usepackage{amssymb}
\usepackage{enumerate}
\usepackage{tikz}
\usepackage{csquotes}
\usepackage{xcolor}
\usepackage[shortlabels]{enumitem}
\usepackage{float}
\newtheorem{theorem}{Theorem}
\newcommand{\del}{{\mathrm{del}}}
\newcommand{\lk}{{\mathrm{lk}}}
\title{On a complete characterization of path-free complexes associated with complete multipartite graphs}
\author{Priyavrat Deshpande, Shuchita Goyal, Rutuja Sawant}
\date{Source: arXiv:2607.05358v1 (CC BY 4.0)}
\begin{document}
\maketitle

\noindent\emph{Source and licence.} On a complete characterization of path-free complexes associated with complete multipartite graphs. Version arXiv:2607.05358v1 (CC BY 4.0), distributed by arXiv under the Creative Commons Attribution 4.0 International licence, http://creativecommons.org/licenses/by/4.0/, as recorded in the arXiv licence metadata for this exact version and shown on the abstract page https://arxiv.org/abs/2607.05358v1. Full licence notice: \texttt{licenses/arxiv-2607.05358-CC-BY-4.0.txt}.\par\smallskip

\noindent\textit{Editorial scope.} Counted unit: the article's theorem V.D (source0.tex lines 946--962). The article's complete original proof of it is delivered with it (source0.tex lines 964--1054, one \texttt{proof} environment of 91 lines). No mathematics is written, repaired or re-derived: the statement and the proof are the article's own lines, in the article's own notation, and no retained line is substituted or re-flowed.  No further source-local context is needed: every cross-reference of every retained line resolves inside the two delivered spans.  Nothing was omitted from any retained span, so the package carries no omission record.  No retained line contains a citation, so the wrapper carries no bibliography and no bibliography entry.  No retained line contains a figure environment, an image-inclusion command or drawing code, so no image is delivered and the package declares no asset dependency.  Editorial formatting only, in addition: an AMS \texttt{amsart} preamble that restates the article's own declarations and macros used by the retained lines, namely the environments \texttt{theorem} on separate counters, together with the standard packages those lines need.  The retained environments print in this document as Theorem 1 in order of appearance, whereas the article prints the same items with the numbering of its own full document.  The complete native source of the article as distributed in the arXiv e-print for this exact version is bundled unchanged as \texttt{sources/204537/source0.tex}.
\begin{theorem}\label{them: General V.D}
Let $V$ be a vertex set and let $W \subseteq V$. Suppose that
\[
-1 \le a \le |W^{c}|-1
\qquad \text{and} \qquad
-1 \le b \le |V|-1.
\]
Then the simplicial complex
\[
\Sigma
=
\left(\Delta_{W} * (\Delta_{W^{c}})^{[a]} \right)
\cup
(\Delta_{V})^{[b]}
\]
is vertex decomposable.
\end{theorem}

\begin{proof}
Suppose $b \ge a+|W|$. Every face of $\Delta_W * (\Delta_{W^c})^{[a]}$ has dimension at most $a+|W|$, so it is a face of $(\Delta_V)^{[b]}$. Thus $\Delta_W * (\Delta_{W^c})^{[a]}$ is a subcomplex of $(\Delta_V)^{[b]}$, and hence $\Sigma=(\Delta_V)^{[b]}$. Since the pure skeleton of simplices is vertex decomposable, $\Sigma$ is vertex decomposable.

Thus we may assume that
\[
b<a+|W|.
\]

If $a=|W^c|-1$, then $(\Delta_{W^c})^{[a]}=\Delta_{W^c}$, and hence $\Delta_W * \Delta_{W^c} = \Delta_V$. Similarly, if $b=|V|-1$, then $(\Delta_V)^{[b]}=\Delta_V$. Also, if $W = V$ then $\Sigma = \Delta_V$. In all cases, $\Sigma=\Delta_V$, and therefore $\Sigma$ is vertex decomposable.

Hence we may further assume that $W \subset V$, 
\[
-1\le a<|W^c|-1
\qquad \text{and} \qquad
-1\le b<|V|-1.
\]

We now prove the statement by induction on $|V|$.

\medskip

\noindent
\textbf{Base case:}
If $|V|=1$, then $\Sigma$ is a simplex. Hence $\Sigma$ is vertex decomposable.

\medskip
\noindent
\textbf{Induction hypothesis:}
Assume that the statement holds for all graphs with the number of vertices strictly smaller than $n$.

\medskip

\noindent
\textbf{Induction step:}
Suppose $|V|=n$.

Choose a vertex $v \in W^c$. 
Such a vertex exists since $W \subset V$ implies that $W^c\neq \emptyset$.

We show that $v$ is a shedding vertex.

First consider the deletion:
\[
\del_\Sigma(v)
=
\left(
\Delta_W * (\Delta_{W^c\setminus\{v\}})^{[a]}
\right)
\cup
(\Delta_{V\setminus\{v\}})^{[b]}.
\]

Since $a<|W^c|-1$, we obtain $a\le |W^c\setminus\{v\}|-1$. 
Also, $b\le |V\setminus\{v\}|-1$, because $|V\setminus\{v\}|=|V|-1$.

Therefore $\del_\Sigma(v)$ is of the same form as in the statement, but on the smaller vertex set $V\setminus\{v\}$. 
By the induction hypothesis, the deletion $\del_\Sigma(v)$ is vertex decomposable.

Next consider the link:
\[
\lk_\Sigma(v)
=
\left(
\Delta_W * (\Delta_{W^c\setminus\{v\}})^{[a-1]}
\right)
\cup
(\Delta_{V\setminus\{v\}})^{[b-1]}.
\]

If $a=-1$, then
($(\Delta_{W^c})^{[a]}=\emptyset$, and hence $\lk_\Sigma(v)=(\Delta_{V\setminus\{v\}})^{[b-1]}$.) Otherwise, $-1\le a-1\le |W^c\setminus\{v\}|-1$.

In either case, $\lk_\Sigma(v)$ is of the same form as in the statement on the vertex set $V\setminus\{v\}$. 
Since $-1\le b-1\le |V\setminus\{v\}|-1$, the induction hypothesis implies that $\lk_\Sigma(v)$ is vertex decomposable.

It remains to verify the shedding condition.
Let $F$ be a facet of $\lk_\Sigma(v)$. 
Then either $F=W\cup F'$, where $F'$ is a facet of  $(\Delta_{W^c\setminus\{v\}})^{[a-1]}$, or $F$ is a facet of $(\Delta_{V\setminus\{v\}})^{[b-1]}$

First suppose $F=W\cup F'$. Since $F'$ is a facet of
\[
(\Delta_{W^c\setminus\{v\}})^{[a-1]},
\]
we have $|F'|=a$, so $|F|=|W|+a$. Because $b<a+|W|$, we get $|F|>b$, hence $F$ is not a face of $(\Delta_{V\setminus\{v\}})^{[b]}$.

The inequality $a<|W^c|-1$ yields a vertex $w\in W^c\setminus(F'\cup\{v\})$. Then $F'\cup\{w\}$ is a face of $(\Delta_{W^c\setminus\{v\}})^{[a]}$, so $F\cup\{w\}=W\cup(F'\cup\{w\})$ is a face of $\del_\Sigma(v)$. Thus $F$ is not a facet of $\del_\Sigma(v)$.

Now suppose $F$ is a facet of $(\Delta_{V\setminus\{v\}})^{[b-1]}$. Then $|F|=b$. Since $b<|V|-1$, there exists $w\in V\setminus(F\cup\{v\})$, and $F\cup\{w\}$ is a face of $(\Delta_{V\setminus\{v\}})^{[b]}$, hence of $\del_\Sigma(v)$. Thus $F$ is not a facet of $\del_\Sigma(v)$.

Therefore no facet of $\lk_\Sigma(v)$ is a facet of $\del_\Sigma(v)$, so $v$ is a shedding vertex. Since both $\del_\Sigma(v)$ and $\lk_\Sigma(v)$ are vertex decomposable, $\Sigma$ is vertex decomposable.
\end{proof}

\end{document}
